Do local income taxes cause residential sorting?
Pennsylvania's Earned Income Tax follows a piecewise max-rule where a
resident pays the higher of the home-side rates and the workplace
rate. The rule creates variation that identifies a causal effect. I
decompose statutory increases by whether the home-side stack already
binds, becomes binding, or never binds, and the never-binding cell is a within-design placebo. I estimate a continuous stacked
difference-in-differences design on block-level LEHD/LODES data
aggregated to municipalities over 2002--2019. The average causal response to a binding increase is a $0.437$
log-point decline per percentage point in the count of workers who
both live and work in their own municipality, about a $10\%$ decline
at the average cumulative increase of $0.25$ percentage points, while
the placebo cell is null. A spline specification traces the full
dose-response, the total and marginal effect at each dose: the effect
accumulates steeply across the increases and flattens as they
grow. The largest decline comes from the first increase, and each subsequent
increase moves fewer workers, consistent with a shrinking pool of
residents willing to relocate. The binding cell passes a set of pre-trend tests and shows no
effect at three fake-timing placebo dates. Controlling for municipal, school-district, and county effective
property-tax rates shows the property tax's effect is about a third of
the EIT's per percentage point.
